\section{Discussion}
\label{sec:discussion}

\subsection{Speed versus quality}
The three evaluated models occupy distinct points on the speed--quality frontier.
The CNNs (ours, $0.18$\,s; Zhang17, $0.22$\,s; pretrained, $0.17$\,s) are
single forward passes and run at interactive rates even on a $6$\,GB GPU.
DeOldify, at $0.51$\,s, buys its perceptual edge with a heavier
self-attention U-Net and a larger render resolution. For automatic
colorization under a latency budget, the CNN is the rational default; the GAN
is worth its cost only when perceptual quality is paramount.

\subsection{Saturation and the desaturation trap}
The classification-plus-annealed-mean design (Eqs.~\ref{eq:loss}--\ref{eq:annealed})
exists precisely to avoid the desaturated mean that an $\ell_2$ regressor
produces. It largely works: our outputs are vivid, not sepia. Yet the residual
LPIPS gap to DeOldify shows the effect is incomplete---on large, color-ambiguous
regions our model still hedges toward muted tones, whereas the adversarially-trained
DeOldify commits to a confident, realistic palette. Class rebalancing and the
$T=0.38$ temperature mitigate but do not eliminate the prior pull toward gray.

\subsection{Failure cases}
Two failure modes recur. First, Zhang17 in automatic mode produces saturated
color blotches and bleeding across object boundaries---an artifact of a network
expecting sparse hints it never receives. Second, all CNNs struggle with
semantically ambiguous content (a shirt that could be any color), defaulting to
the most frequent training color; this is fundamental to the task rather than to
any single model.

\subsection{Paradigm-level reading}
\begin{itemize}
  \item \textbf{CNN} --- fast, robust, fully automatic; the strongest
        speed--quality compromise for unattended colorization.
  \item \textbf{Interactive CNN} --- excellent \emph{with} user hints, but a poor
        automatic colorizer; the wrong tool for a hands-off pipeline.
  \item \textbf{GAN} --- highest perceptual quality and best restoration
        behavior, at $\sim$$3\times$ the latency.
\end{itemize}

\subsection{The diffusion paradigm: a reproducibility finding}
The decision to exclude diffusion is itself a finding about the open-source
deep-learning ecosystem rather than a property of the paradigm. Every
well-trained colorization checkpoint we surveyed in 2026 was built on Stable
Diffusion 2.1, and the upstream deprecation of SD~2.x removed the foundation
out from under that downstream collection in one stroke. The lesson generalises:
a paradigm whose practical implementations share a single foundation model is
exposed to that foundation model's lifecycle decisions. A future replication
should revisit the diffusion slot once a maintained colorization checkpoint
exists on a still-supported backbone---neither of which we could identify at
the time of writing.

\subsection{Limitations}
The ``pretrained'' row is not a faithful ECCV baseline (head mismatch), so we
avoid drawing fine-tuning-gain conclusions from it. The benchmark is $1{,}000$
images on a single GPU with no human preference study; perceptual quality is
proxied by LPIPS rather than measured directly. Finally, the comparison is
three paradigms only---the diffusion paradigm is an explicit gap, discussed
above, rather than a measured frontier point.
